% NeurIPS style preprint, recreated with standard packages only.
\documentclass[10pt]{article}
\usepackage[letterpaper,textwidth=5.5in,textheight=9in,top=1in]{geometry}
\usepackage{mathptmx}
\usepackage{amsmath}
\usepackage{amssymb}
\usepackage{titlesec}

\titleformat{\section}{\large\bfseries\sffamily}{\thesection}{0.7em}{}
\titleformat{\subsection}{\normalsize\bfseries\sffamily}{\thesubsection}{0.7em}{}
\titlespacing*{\section}{0pt}{1.6ex plus .4ex}{0.9ex}
\titlespacing*{\subsection}{0pt}{1.3ex plus .3ex}{0.7ex}

\begin{document}

\noindent\rule{\textwidth}{1.2pt}
\begin{center}
  \vspace{4pt}
  {\sffamily\bfseries\LARGE Gradient Sketching for Communication-Efficient\\ Federated Learning\par}
  \vspace{6pt}
\end{center}
\noindent\rule{\textwidth}{0.5pt}

\vspace{14pt}
\begin{center}
\begin{minipage}[t]{0.31\textwidth}\centering
  {\bfseries Maya Lindqvist}\\ Department of Computer Science\\ Aalto University\\ \texttt{maya.lindqvist@aalto.fi}
\end{minipage}\hfill
\begin{minipage}[t]{0.31\textwidth}\centering
  {\bfseries Daniel Okafor}\\ School of Informatics\\ University of Edinburgh\\ \texttt{d.okafor@ed.ac.uk}
\end{minipage}\hfill
\begin{minipage}[t]{0.31\textwidth}\centering
  {\bfseries Elena Marchetti}\\ Institute for ML Systems\\ ETH Zurich\\ \texttt{elenam@ethz.ch}
\end{minipage}
\end{center}

\let\thefootnote\relax\footnotetext{Preprint. Under review.}

\vspace{16pt}
\begin{center}{\sffamily\bfseries\large Abstract}\end{center}
\begin{list}{}{\leftmargin=0.4in\rightmargin=0.4in}\item\relax
Federated optimization is dominated by the cost of shipping dense gradient
updates between clients and the server. We propose gradient sketching, a
compression scheme that projects each client update onto a shared random
subspace and reconstructs an unbiased estimate on the server. The sketch
dimension adapts to the observed gradient spectrum, so communication shrinks
automatically as training converges. On image and language benchmarks with up
to 1{,}024 simulated clients, our method cuts uplink traffic by 14x while
matching the accuracy of uncompressed FedAvg within 0.3 points.
\end{list}

\section{Introduction}
Cross-device federated learning trains a shared model on data that never
leaves user hardware. In practice the bottleneck is not computation but the
uplink: mobile clients upload millions of parameters per round over
constrained links. Prior work reduces this cost with quantization or top-$k$
sparsification~\cite{stromberg2023topk}, but both approaches introduce biased
updates that require error feedback to remain stable.

We take a different route. Each round, the server broadcasts a seed that
defines a random projection $S_t \in \mathbb{R}^{k \times d}$ with $k \ll d$.
Clients upload the sketch $S_t g_i$ of their local gradient $g_i$, and the
server forms the unbiased reconstruction
\begin{equation}
  \hat{g}_t \;=\; \frac{d}{k}\, S_t^{\top} \Big( \tfrac{1}{n} \sum_{i=1}^{n} S_t\, g_i \Big),
  \label{eq:sketch}
\end{equation}
which it feeds to any server-side optimizer. Because the projection is shared,
sketches aggregate before reconstruction and the server cost is a single
matrix product.

\section{Method and Analysis}
Equation~\eqref{eq:sketch} yields an estimator whose variance is controlled by
the stable rank of the averaged gradient, which we observe to fall sharply
after the first few epochs. We exploit this by shrinking $k$ on a schedule
tied to a cheap online estimate of the spectrum, following the adaptive
compression analysis of~\cite{vasquez2024adaptive}. Under standard smoothness
assumptions the scheme retains the $O(1/\sqrt{T})$ rate of FedAvg, and unlike
sign-based compressors it needs no error accumulation
buffers~\cite{herzog2022signfl}. A full proof and ablations over sketch
schedules appear in the appendix.

\begin{thebibliography}{9}
\bibitem{stromberg2023topk} L. Str\"omberg and P. Aravind.
Top-$k$ sparsification revisited: error feedback in cross-device settings.
In \textit{Proceedings of the 40th International Conference on Machine Learning}, 2023.
\bibitem{vasquez2024adaptive} R. V\'asquez, T. Nakayama, and S. Bhat.
Adaptive compression rates for distributed SGD.
\textit{Journal of Machine Learning Research}, 25(118):1--43, 2024.
\bibitem{herzog2022signfl} K. Herzog and A. Deml.
SignFL: one-bit federated optimization at scale.
In \textit{Advances in Neural Information Processing Systems 35}, 2022.
\end{thebibliography}

\end{document}
